\documentclass{../common/latex/delta_base_styles/delta_base_styles}

\addbibresource{example_refs.bib}

% Safe imports
\subimport{../common/latex/delta_header/}{delta_header.tex}
\subimport{../common/latex/equipment_table/}{equipment_table.tex}
\subimport{../common/latex/signature_table/}{signature_table.tex}
\subimport{../common/latex/version_approval_table/}{version_approval_table.tex}
\graphicspath{
  {../common/latex/delta_base_styles/}
  {../common/latex/delta_header/}
  {../common/latex/equipment_table/}
  {../common/latex/signature_table/}
  {../common/latex/version_approval_table/}
}

% Set metadata
\doctitle{MBSE-First Verification Report Template}
\docsubtitle{Report-by-reference for Capella/ARCADIA IADT verification activities}

\docid{TEST-000}
\version{1.0}
\docdate{\today}
\preparedby{Alice}
\approvedby{Bob}
\observations{Draft template}

\newcommand{\placeholder}[1]{\textcolor{DarkCyan}{#1}}

\begin{document}

\makedeltaheader

\section{Report Control and Verification Activity Summary}
\makeversionapprovaltable

\placeholder{Use this template after the verification activity has been researched in the subsystem \texttt{tests/README.md} and modeled in Capella/D2 under \texttt{tests/<activity-id>/}. The Capella/ARCADIA model is the source of truth. This report records the executed evidence, configuration, results, deviations, and pass/fail rationale; it should not duplicate stable model content unnecessarily.}

\renewcommand{\arraystretch}{1.3}
\setlength{\tabcolsep}{3pt}
\begin{tabular}{|>{\centering\arraybackslash}m{3.1cm}
                |>{\centering\arraybackslash}m{4.1cm}
                |>{\centering\arraybackslash}m{3.6cm}
                |>{\centering\arraybackslash}m{4.1cm}|}
\hline
\textbf{Field} & \textbf{Value} & \textbf{Field} & \textbf{Value} \\
\hline
Verification activity ID & \placeholder{e.g., DPS-V10-T-001} & Report type & \placeholder{Activity report / campaign report / retest report} \\
\hline
Subsystem / SSIV & \placeholder{Subsystem and development version, e.g., PDS \& ESS v1.0} & IADT method & \placeholder{Inspection / Analysis / Demonstration / Testing} \\
\hline
IVV source category & \placeholder{component/link, component exchange, allocation, functional chain/scenario, or constraint} & Traceability target(s) & \placeholder{mission, capability, use case, feared event, or constraint IDs} \\
\hline
Model baseline & \placeholder{Capella/D2 commit, tag, branch, or controlled baseline ID} & As-verified configuration & \placeholder{HW revision, firmware commit, article serial numbers, tools} \\
\hline
Modeled definition folder & \placeholder{\texttt{tests/<activity-id>/}} & Evidence folder & \placeholder{\texttt{tests/results/<activity-id>/}} \\
\hline
Overall result & \placeholder{Pass / Fail / Conditional pass / Inconclusive / Not run} & Retest or waiver reference & \placeholder{N/A or identifier} \\
\hline
\end{tabular}

\section{Purpose, Scope, and MBSE Source of Truth}
\subsection{Purpose}
\placeholder{State why this verification was executed and which modeled verification activity it closes or partially closes.}

\subsection{Scope}
\begin{itemize}
  \item \textbf{Included model elements:} \placeholder{List the referenced physical components, physical links, component exchanges, logical components/functions, functional chains/scenarios, or constraints by model name/ID.}
  \item \textbf{Excluded model elements:} \placeholder{List explicit exclusions and their rationale, such as another SSIV, another subsystem, or a candidate activity not yet modeled.}
  \item \textbf{Applicable SSIV/development version:} \placeholder{State the exact version. If inferred from folders rather than a controlled SSIV, state that limitation.}
\end{itemize}

\subsection{MBSE-first reporting rule}
\placeholder{If the unit under verification, environment, test means, interfaces, stimuli, expected behavior, and pass/fail constraints are already defined in the model, reference them here instead of repeating them. Add details only when needed to identify the as-tested configuration, actual execution conditions, observed results, or deviations from the modeled definition.}

\section{Referenced MBSE Definition and Traceability}
\subsection{Source and Verification Definition Views}
\begin{tabular}{|>{\centering\arraybackslash}m{3.1cm}
                |>{\centering\arraybackslash}m{5.0cm}
                |>{\centering\arraybackslash}m{7.0cm}|}
\hline
\textbf{Reference type} & \textbf{Path or model element} & \textbf{Notes / baseline} \\
\hline
Subsystem test plan & \placeholder{\texttt{<subsystem>/MBSE/tests/README.md}} & \placeholder{Planning row/status for this activity.} \\
\hline
Research references & \placeholder{\texttt{<subsystem>/MBSE/tests/references/}} & \placeholder{Standards, literature, statistics, and fault-hardening references used during planning.} \\
\hline
Source architecture views & \placeholder{\texttt{<subsystem>/MBSE/<version>/...}} & \placeholder{Unmodified Capella/D2 physical architecture baseline.} \\
\hline
Modeled verification views & \placeholder{\texttt{<subsystem>/MBSE/tests/<activity-id>/}} & \placeholder{Copied and extended views with verification actors, equipment, functions, links, chains, constraints, and traceability.} \\
\hline
Functional chains / scenarios & \placeholder{Model element names or diagram paths} & \placeholder{Only for demonstration/testing or chain-based analyses.} \\
\hline
Pass/fail constraints & \placeholder{Constraint IDs/names} & \placeholder{Where each acceptance criterion is modeled.} \\
\hline
\end{tabular}

\subsection{Traceability and Coverage Summary}
\begin{tabular}{|>{\centering\arraybackslash}m{2.3cm}
                |>{\centering\arraybackslash}m{3.0cm}
                |>{\centering\arraybackslash}m{3.1cm}
                |>{\centering\arraybackslash}m{3.2cm}
                |>{\centering\arraybackslash}m{3.2cm}|}
\hline
\textbf{Activity ID} & \textbf{Model elements covered} & \textbf{Trace target} & \textbf{Viewpoints} & \textbf{Result / coverage impact} \\
\hline
\placeholder{ID} & \placeholder{PC/PL, CE, allocation, chain, scenario, or constraint} & \placeholder{mission/capability/use case/feared event/constraint} & \placeholder{statistics, fault hardening, environment, safety, timing, etc.} & \placeholder{Closed / partially closed / gap remains} \\
\hline
\end{tabular}

\section{References}
\placeholder{List documents, standards, modeled references, and procedures used to research, define, or execute this activity. The planning prompts are process references; the controlled verification content is the resulting model and test plan.}
\begin{itemize}
  \item \textbf{Project IVV policy:} \placeholder{\texttt{PM\&SE/IVV.md}.}
  \item \textbf{Research workflow:} \placeholder{\texttt{/Users/antho/agent-tools/syseng/prompts/capella-pa-tests-research.md}.}
  \item \textbf{Model-definition workflow:} \placeholder{\texttt{/Users/antho/agent-tools/syseng/prompts/capella-pa-tests-definition.md}.}
  \item \textbf{Subsystem plan and modeled definition:} \placeholder{\texttt{tests/README.md} row and \texttt{tests/<activity-id>/} diagrams.}
  \item \textbf{Standards:}
  \nocite{ISO29119,IEC60068,ISO17025}
  \printbibliography[heading=none]
  \item \textbf{Additional saved references:} \placeholder{List local files in \texttt{tests/references/} used for statistics, environmental assumptions, fault hardening, datasheet limits, or workmanship criteria.}
\end{itemize}

\section{Definitions and Abbreviations}
\begin{itemize}
  \item \textbf{ARCADIA:} Architecture Analysis and Design Integrated Approach.
  \item \textbf{Capella/D2 view:} Modeled system or verification definition used as the authoritative verification source.
  \item \textbf{IADT:} Inspection, Analysis, Demonstration, and Testing.
  \item \textbf{MBSE:} Model-Based Systems Engineering.
  \item \textbf{SSIV:} System or Subsystem Integration Version.
  \item \textbf{UUT:} Unit Under Verification; use DUT/EUT only if required by a referenced standard.
  \item \textbf{Deviation:} Any executed condition, setup, criterion, or sequence that differs from the modeled verification definition.
\end{itemize}

\section{Test Item and As-Verified Configuration}
\placeholder{Identify only execution-specific configuration details that are not already controlled by the model, or that must be bound to evidence for repeatability.}

\begin{tabular}{|>{\centering\arraybackslash}m{2.7cm}
                |>{\centering\arraybackslash}m{2.7cm}
                |>{\centering\arraybackslash}m{2.7cm}
                |>{\centering\arraybackslash}m{3.1cm}
                |>{\centering\arraybackslash}m{3.4cm}|}
\hline
\textbf{Item / article} & \textbf{Serial or revision} & \textbf{Manufacturer / owner} & \textbf{HW/SW configuration} & \textbf{Photographs / diagrams / evidence} \\
\hline
\placeholder{UUT or analyzed artifact} & \placeholder{Serial, board rev, firmware commit} & \placeholder{Name} & \placeholder{As-tested configuration} & \placeholder{Photo/log/model reference} \\
\hline
\end{tabular}

\section{Verification Means, Environment, Equipment, and Operators}
\subsection{Modeled Verification Means}
\placeholder{Reference the external actors, test equipment, human operators, verification-only functions, physical links, component exchanges, and environmental constraints added in the modeled definition.}

\subsection{Actual Environmental Conditions}
\begin{itemize}
  \item \textbf{Date/time/location:} \placeholder{Actual execution context.}
  \item \textbf{Temperature, humidity, pressure, EMC/EMI, lighting, ESD, safety controls:} \placeholder{Record values when relevant to the modeled criteria or evidence interpretation.}
  \item \textbf{Configuration stability:} \placeholder{State warm-up, calibration status, battery state, radio settings, or other conditions that affect the result.}
\end{itemize}

\subsection{Equipment Used}
\equipmenttable{
  1 & \placeholder{Equipment name} & \placeholder{Type or model} & \placeholder{Manufacturer} & \placeholder{Calibration date / N.A.} & \placeholder{Serial no.} & \placeholder{Role in modeled definition} \\
  \hline
}

\subsection{Operators and Roles}
\begin{tabular}{|>{\centering\arraybackslash}m{3.3cm}
                |>{\centering\arraybackslash}m{4.2cm}
                |>{\centering\arraybackslash}m{4.2cm}
                |>{\centering\arraybackslash}m{3.2cm}|}
\hline
\textbf{Name} & \textbf{Modeled actor/function} & \textbf{Actual role performed} & \textbf{Training / authorization} \\
\hline
\placeholder{Operator} & \placeholder{e.g., press launch button, record measurement, inspect PCB} & \placeholder{Actual role} & \placeholder{Reference or N/A} \\
\hline
\end{tabular}

\section{As-Executed Procedure or Analysis Method}
\placeholder{Reference the modeled functional chain, scenario, inspection checklist, or analysis checklist. Record the executed sequence or analysis method only to the level needed to reproduce evidence and identify deviations.}

\begin{tabular}{|>{\centering\arraybackslash}m{0.9cm}
                |>{\centering\arraybackslash}m{3.2cm}
                |>{\centering\arraybackslash}m{3.2cm}
                |>{\centering\arraybackslash}m{3.2cm}
                |>{\centering\arraybackslash}m{3.2cm}|}
\hline
\textbf{Step} & \textbf{Modeled action / analysis item} & \textbf{Actual execution / evidence} & \textbf{Deviation from model} & \textbf{Initial result} \\
\hline
1 & \placeholder{Setup/action/check} & \placeholder{Log/photo/script/checklist row} & \placeholder{None or ID} & \placeholder{Pass/Fail/Obs.} \\
\hline
2 & & & & \\
\hline
3 & & & & \\
\hline
\end{tabular}

\section{Results and Evidence Index}
\subsection{Evidence Inventory}
\begin{tabular}{|>{\centering\arraybackslash}m{2.4cm}
                |>{\centering\arraybackslash}m{4.0cm}
                |>{\centering\arraybackslash}m{3.0cm}
                |>{\centering\arraybackslash}m{5.0cm}|}
\hline
\textbf{Evidence ID} & \textbf{File/path or artifact} & \textbf{Produced by} & \textbf{Purpose / linked model element} \\
\hline
\placeholder{E-001} & \placeholder{\texttt{results/<activity-id>/...}} & \placeholder{tool/operator} & \placeholder{raw data, photo, script output, checklist, analysis note} \\
\hline
\end{tabular}

\subsection{Raw Results}
\placeholder{Summarize raw measured values, observations, checklist counts, analysis findings, logs, screenshots, photos, or scripts. Keep raw artifacts in the evidence folder and cite them here.}

\subsection{Statistical Treatment and Uncertainty}
\placeholder{Apply the statistical policy from \texttt{PM\&SE/IVV.md} and the modeled activity: exact binomial confidence for success/failure or PDR claims, at least 30 samples for continuous measurements where applicable, 59 in-limit samples for 95/95 deadline claims, or state why statistics are not applicable for inspection/analysis. Record confidence, sample size, independence assumptions, uncertainty, and limitations.}

\subsection{Fault-Hardening Results}
\placeholder{Record injected or analyzed fault cases, forbidden events, recovery behavior, and whether failures were detected, rejected, bounded, logged, or safely handled as modeled.}

\section{IADT-Specific Pass/Fail Rationale}
\placeholder{Complete the subsection(s) that apply to this activity. Mark non-applicable subsections as N/A with rationale.}

\subsection{Physical Component and Physical Link Inspection}
\placeholder{For component/link inspection activities, state the inspected PC/PL population, inspection coverage, missing/extra items, continuity/polarity/short results, workmanship findings, and disposition. Pass only if all modeled physical components and links are present or explicitly dispositioned.}

\subsection{Component Exchange Analysis}
\placeholder{For component-exchange analyses, state each analyzed CE, evidence that the exchange exists, protocol/direction/data semantics, implementation consistency with model intent, and fault-hardening rationale.}

\subsection{Functional Allocation Analysis}
\placeholder{For allocation analyses, state whether every intended function is allocated to the modeled logical component, no foreign functions are allocated, and cross-boundary flows use modeled component exchanges or documented dispositions.}

\subsection{Functional Chain or Scenario Testing/Demonstration}
\placeholder{For chain/scenario activities, state stimuli, expected behavior, observed behavior, pass/fail constraints, trace target, relevant viewpoints, sample counts, timing/rate/accuracy metrics, and evidence that the human/equipment verification means participated as modeled.}

\subsection{Constraint-Derived Verification}
\placeholder{For constraint-derived activities, state each constraint, where it is modeled as a pass/fail condition, the measured/analyzed result, and whether the constraint is closed, partially closed, or still requires another verification activity.}

\section{Deviations, Anomalies, Waivers, Assumptions, and Limitations}
\begin{tabular}{|>{\centering\arraybackslash}m{2.0cm}
                |>{\centering\arraybackslash}m{2.4cm}
                |>{\centering\arraybackslash}m{5.2cm}
                |>{\centering\arraybackslash}m{3.1cm}
                |>{\centering\arraybackslash}m{2.6cm}|}
\hline
\textbf{ID} & \textbf{Type} & \textbf{Description and impact} & \textbf{Disposition / owner} & \textbf{Status} \\
\hline
\placeholder{DEV-001} & \placeholder{Deviation / anomaly / waiver / assumption / limitation} & \placeholder{Describe impact on interpretation of result.} & \placeholder{Accept, retest, update model, open issue} & \placeholder{Open/Closed} \\
\hline
\end{tabular}

\section{Coverage and SSIV Completeness Impact}
\placeholder{Update the subsystem \texttt{tests/README.md} after report approval if this result changes status, coverage, gaps, or expected report references.}

\begin{itemize}
  \item \textbf{Coverage closed by this report:} \placeholder{Activities/model elements/constraints now evidenced.}
  \item \textbf{Coverage partially closed:} \placeholder{Items needing additional evidence, model update, or retest.}
  \item \textbf{Remaining gaps:} \placeholder{Missing traceability, missing modeled CE/allocation/constraint, unconfirmed SSIV, insufficient statistics, missing fault case, or missing evidence.}
  \item \textbf{Impact on SSIV readiness:} \placeholder{Ready / ready with limitations / not ready, with rationale.}
\end{itemize}

\section{Conclusion}
\begin{itemize}
  \item \textbf{Overall result:} \placeholder{Pass / Fail / Conditional pass / Inconclusive / Not run.}
  \item \textbf{Pass/fail rationale:} \placeholder{Tie the result to modeled pass/fail constraints and actual evidence.}
  \item \textbf{Required follow-up:} \placeholder{Retest, update Capella/D2 model, update tests/README.md, raise anomaly, or accept waiver.}
\end{itemize}

\section{Signatures}

\makeatletter
\signatureTable{
\@approvedby & Approver & \@docdate & \\
\hline
}
\makeatother

\section{Annexes}
\placeholder{Attach or reference controlled supporting material:}
\begin{itemize}
  \item Modeled verification diagrams exported from \texttt{tests/<activity-id>/}.
  \item Raw data, scripts, notebooks, logs, photos, screenshots, videos, and checklist scans.
  \item Calibration certificates, datasheets, environmental profiles, safety documentation, and workmanship records.
  \item Updated coverage table excerpt from the subsystem \texttt{tests/README.md}, if needed.
\end{itemize}

\vspace{2cm}
\begin{center}
\Huge \textcolor{DarkCyan}{\textbf{IMPORTANT:} \\[0.5cm]}
\large
\textcolor{DarkCyan}{This is an MBSE-first report template. Replace all dark-cyan guidance with actual evidence. If a verification activity is not yet represented in the Capella/D2 model or the subsystem tests plan, stop and define it first using the project MBSE verification workflow. The final report should reference the model definition and record the as-executed evidence, not become the uncontrolled source of the test design.}
\end{center}

\end{document}
